% Generated draft: requires mathematical and source review.
\documentclass[11pt]{article}
\usepackage[margin=25mm]{geometry}
\usepackage{fontspec}
\setmainfont{Latin Modern Roman}
\usepackage{amsmath,amssymb,mathtools,unicode-math}
\setmathfont{Latin Modern Math}
\usepackage{enumitem,needspace,xurl,hyperref}
\setlength{\parindent}{0pt}
\setlength{\parskip}{5pt}
\setlength{\emergencystretch}{3em}
\newcommand{\N}{\mathbb{N}}
\newcommand{\Z}{\mathbb{Z}}
\newcommand{\Q}{\mathbb{Q}}
\newcommand{\R}{\mathbb{R}}
\newcommand{\C}{\mathbb{C}}
\newcommand{\F}{\mathbb{F}}
\newcommand{\A}{\mathbb{A}}
\newcommand{\PP}{\mathbb{P}}
\newcommand{\E}{\mathbb{E}}
\newcommand{\eps}{\varepsilon}
\newcommand{\dd}{\,\mathrm{d}}
\begin{document}
\section*{UnsolvedMath ID: 30000695}
\textbf{Source problem code:} catalogue entry
\subsection{Definitions and mathematical statement}
In the parabolic Anderson model, determine the geometry and statistics of the ``relevant islands'' governing long-time behavior. A standard static-potential formulation is
\[
 \partial_tu(t,x)=\kappa\Delta u(t,x)+\xi(x)u(t,x),
 \qquad u(0,x)=u_0(x),
\]
where $\xi$ is a random potential. This equation explains the terminology; the source does not specify discrete or continuum space, its Laplacian, dimension, diffusion constant, potential law and dependence, or initial condition.

The original asks all three of the following:
\begin{enumerate}
\item Determine how large the relevant islands are, and how large both the potential and the solution are within them. This includes spatial island scales and potential/solution magnitude scales, not just one of these quantities.
\item Determine the shapes of both the potential and the solution in the relevant islands, with the necessary centering and rescaling.
\item Determine how many relevant islands there are and their spatial distribution.
\end{enumerate}

A precise formulation must define which regions count as relevant islands and which observable they dominate, for example by a specified mass-concentration criterion rather than a chosen informal notion of high potential. The source gives no such criterion, observation window, normalization, shape topology, counting rule, or mode of probabilistic convergence. It also does not distinguish quenched assertions for fixed environments from annealed assertions averaged over the potential. The count on an infinite state space is not silently identified with the count in a finite growing window.

The source describes the relevant-island picture as its long-time context; the draft does not prove localization for an unspecified potential law. The supplied assessment reports known island shapes and localization in important tail classes while saying the general size, shape, number, and distribution questions remain model dependent. No universal scale, profile, point process, or tail assumption is inserted.
\subsection{Short English statement}
The question asks for the size and heights, potential and solution profiles, and number and locations of the islands controlling parabolic-Anderson long-time behavior. The potential law, initial data, relevance criterion, scales, and probabilistic mode are missing.
\subsection{Sources}
\begin{itemize}
\item Local UnsolvedMath registry, permanent ID 30000695, code catalogue entry, ``Relevant Islands in the Parabolic Anderson Model.'' Proposed year 2007; proposer unspecified.
\item \emph{Mini-Workshop: Multiscale and Variational Methods in Material Science and Quantum Theory of Solids} (2007), Oberwolfach Reports, pp.\ 371--416, DOI 10.4171/owr/2007/06; \url{https://doi.org/10.4171/owr/2007/06}.
\item Supplied reference ``geometric intermittency'': \url{https://arxiv.org/abs/math/0507585}.
\item Supplied ``parabolic Anderson model monograph'': \url{https://link.springer.com/book/10.1007/978-3-319-33596-4}.
\item Statement status exact; the supplied verification calls all three localization questions clear and complete and says they were retained byte-for-byte. The missing model and asymptotic conventions still require recovery.
\item Literature triage checked 2026-08-21: open. It reports results for important tail classes but describes the general island questions as model dependent and unresolved. Exact hypotheses and current status require independent verification. Created 2026-08-25; updated 2026-09-16.
\item Registry curation by Ulam AI / UnsolvedMath Contributors, CC BY 4.0; underlying report and literature provenance remain as supplied. Unreviewed statement-only draft.
\end{itemize}
\end{document}
